\documentclass[11pt]{article}
\usepackage[T1]{fontenc}
\usepackage{amsmath,amssymb,amsthm,mathtools}
\usepackage{graphicx}
\usepackage{xcolor}
\usepackage[margin=1in]{geometry}
\usepackage{microtype}
\usepackage[colorlinks=true,linkcolor=blue!55!black,citecolor=blue!55!black,urlcolor=blue!55!black]{hyperref}

\newtheorem{theorem}{Theorem}[section]
\newtheorem{lemma}[theorem]{Lemma}
\newtheorem{remark}[theorem]{Remark}
\newcommand{\F}{\mathcal F}
\newcommand{\Lip}{\operatorname{Lip}}
\newcommand{\id}{\operatorname{id}}
\newcommand{\R}{\mathbb R}

\title{Open convex domains have the same bidual projection bound\\
\large A candidate full solution to Question 1.5 of arXiv:1703.08384}
\author{Automatic Mathematics Project --- run \texttt{fa\_banach\_001}, lane 3}
\date{9 August 2026}

\begin{document}
\maketitle

\begin{center}
\fbox{\parbox{0.92\linewidth}{
\textbf{Status.} Candidate full solution, likely valid, pending specialist
review.  The answer to Question 1.5 is affirmative, with exactly the
Banach--Mazur bound from Theorem 1.1 of the source paper.
}}
\end{center}

\begin{abstract}
Let $E$ be a finite-dimensional normed space and let $\Omega\subset E$ be
nonempty, open, and convex.  We prove that
\(\F(\Omega)\) is complemented in its bidual by a projection of norm at most
\(d_{BM}(E,\ell_2^{\dim E})\).  The argument first handles the Euclidean case:
the nearest-point map onto $\overline\Omega$ is a nonexpansive retraction, so
its linearization transfers the norm-one bidual projection for the whole
Euclidean space to the domain.  Since $\Omega$ is dense in its closure, their
Lipschitz-free spaces are canonically isometric.  A Banach--Mazur change of
norm then gives the asserted sharp inherited bound.
\end{abstract}

\section{Source question and exact answer}

C\'uth, Kalenda, and Kaplick\'y prove that if $E$ is a real normed space of
finite dimension $d\geq 2$, then there is a projection
\[
 Q_E:\F(E)^{**}\longrightarrow \F(E),\qquad
 \|Q_E\|\leq d_{BM}(E,\ell_2^d).
\]
Their Question 1.5 asks whether the same statement holds for
$\F(\Omega)$ when $\Omega$ is a nonempty open convex subset of $E$
\cite{CKK}.  The source excerpt is reproduced in Figure~\ref{fig:question}.

\begin{figure}[ht]
  \centering
  \includegraphics[width=\linewidth]{figures/open_problem_crop.png}
  \caption{Corollary 1.3 and Questions 1.4--1.6 on page 2 of the source
  paper.  Question 1.5 is the target of this packet.}
  \label{fig:question}
\end{figure}

All Lipschitz-free spaces below are over real scalars.  If $o$ is the
distinguished point of a metric space $M$, we write $\F(M,o)$ when the
basepoint needs to be visible.  Let
\(\kappa_X:X\to X^{**}\) denote the canonical embedding.

\begin{theorem}[Affirmative answer to Question 1.5]\label{thm:main}
Let $(E,\|\cdot\|_E)$ be a real normed space of finite dimension $d\geq 2$,
let $\Omega\subset E$ be nonempty, open, and convex, and choose any
$o\in\Omega$ as basepoint.  There is a bounded linear map
\[
 P_\Omega:\F(\Omega,o)^{**}\longrightarrow \F(\Omega,o)
\]
such that
\[
 P_\Omega\kappa_{\F(\Omega,o)}=\id_{\F(\Omega,o)}
 \quad\text{and}\quad
 \|P_\Omega\|\leq d_{BM}(E,\ell_2^d).
\]
Thus the canonical copy of $\F(\Omega,o)$ is complemented in its bidual
with exactly the estimate in Theorem 1.1 of the source paper.
\end{theorem}

\section{Proof intuition}

There are only two ingredients.  First put a Euclidean norm on the ambient
vector space.  For a closed convex subset $C$ of a finite-dimensional Hilbert
space $H$, the nearest-point map $r:H\to C$ is $1$-Lipschitz and fixes $C$.
The universal property of Lipschitz-free spaces turns $r$ into a contractive
linear retraction $\F(H)\to\F(C)$.  Sandwiching the norm-one projection
$\F(H)^{**}\to\F(H)$ from Corollary 1.3 of the source paper between this
retraction and the bidual of the inclusion yields a norm-one projection onto
$\F(C)$.  Taking $C=\overline\Omega$ costs nothing because a metric space and
its completion have the same Lipschitz-free space.

Second, choose an isomorphism from $E$ to $\ell_2^d$ that realizes the
Banach--Mazur distance.  Its affine restriction to $\Omega$ linearizes to an
isomorphism of free spaces.  Conjugating the Euclidean projection by this
isomorphism multiplies its norm by at most
$\|T\|\|T^{-1}\|=d_{BM}(E,\ell_2^d)$.

\section{The Euclidean-domain step}

We record the elementary functorial facts so that the norm estimate is fully
auditable.

\begin{lemma}[Passing to the completion]\label{lem:completion}
If a pointed metric space $(M,o)$ is dense in a pointed metric space
$(C,o)$, then the inclusion induces an onto linear isometry
\[
 J:\F(M,o)\longrightarrow \F(C,o).
\]
\end{lemma}

\begin{proof}
The linearization $J$ of the inclusion is an isometry.  Indeed, every
real-valued Lipschitz function on $M$ has a unique extension to $C$ with the
same Lipschitz constant, so the dual formula for the free-space norm gives
norm equality on the span of the evaluation vectors and hence on its
closure.  The range of $J$ is closed.  If $x\in C$ and $x_n\in M$ tends to
$x$, then
\[
 \|\delta(x_n)-\delta(x)\|_{\F(C)}=d(x_n,x)\longrightarrow 0.
\]
Thus the range contains every $\delta(x)$, $x\in C$, and is therefore all of
$\F(C,o)$.
\end{proof}

\begin{lemma}[A retract inherits the bidual projection]\label{lem:retract}
Let $(C,o)$ be a pointed metric subspace of $(M,o)$.  Suppose that there is a
$1$-Lipschitz retraction $r:M\to C$.  If
$Q_M:\F(M,o)^{**}\to\F(M,o)$ satisfies
$Q_M\kappa_{\F(M,o)}=\id$ and $\|Q_M\|\leq\lambda$, then there is a map
$Q_C:\F(C,o)^{**}\to\F(C,o)$ satisfying
$Q_C\kappa_{\F(C,o)}=\id$ and $\|Q_C\|\leq\lambda$.
\end{lemma}

\begin{proof}
Let $I:\F(C,o)\to\F(M,o)$ be induced by inclusion and let
$R:\F(M,o)\to\F(C,o)$ be induced by $r$.  McShane extension shows that $I$
is an isometry; the universal property gives $\|R\|\leq1$ and $RI=\id$.
Set
\[
 Q_C=R\,Q_M\,I^{**}.
\]
Then $\|Q_C\|\leq\lambda$.  Naturality of the canonical embeddings gives,
for every $u\in\F(C,o)$,
\[
 Q_C\kappa_{\F(C,o)}u
 =R Q_M\kappa_{\F(M,o)}Iu=RIu=u.
\]
\end{proof}

\begin{lemma}[Euclidean convex domains]\label{lem:euclidean}
Let $H$ be a finite-dimensional Euclidean space and let
$\Omega\subset H$ be nonempty, open, and convex.  For every basepoint
$o\in\Omega$, the space $\F(\Omega,o)$ is $1$-complemented in its bidual.
\end{lemma}

\begin{proof}
Put $C=\overline\Omega$.  This is a nonempty closed convex subset of $H$.
Its metric projection
\[
 r(x)=\operatorname*{argmin}_{c\in C}\|x-c\|_2
\]
is well-defined, fixes $C$, and is nonexpansive.  For completeness, if
$p=r(x)$ and $q=r(y)$, the variational inequalities for the two minimizers
give
\[
 \langle x-p,q-p\rangle\leq0,\qquad
 \langle y-q,p-q\rangle\leq0.
\]
Adding them yields
\[
 \|p-q\|_2^2\leq\langle p-q,x-y\rangle
 \leq\|p-q\|_2\|x-y\|_2,
\]
so $\|r(x)-r(y)\|_2\leq\|x-y\|_2$.  Also $r(o)=o$.

Corollary 1.3 of the source paper supplies a map
$Q_H:\F(H,o)^{**}\to\F(H,o)$ of norm $1$ that is the identity on the
canonical copy of $\F(H,o)$ \cite{CKK}.  Lemma~\ref{lem:retract} applied to
$C\subset H$ and $r$ gives a norm-one bidual projection for $\F(C,o)$.
Finally, $\Omega$ is dense in $C$, so Lemma~\ref{lem:completion} identifies
$\F(\Omega,o)$ isometrically with $\F(C,o)$.  Transporting the projection
through this isometry completes the proof.
\end{proof}

\section{Banach--Mazur transport and proof of the theorem}

\begin{proof}[Proof of Theorem~\ref{thm:main}]
The finite-dimensional Banach--Mazur infimum is attained.  Choose a linear
isomorphism
\[
 T:E\longrightarrow\ell_2^d,
 \qquad \|T\|\,\|T^{-1}\|=d_{BM}(E,\ell_2^d).
\]
To keep basepoints explicit, use the affine map
\[
 A:\Omega\longrightarrow\Omega':=T(\Omega-o),
 \qquad A(x)=T(x-o).
\]
It sends $o$ to $0$, and $\Omega'$ is again nonempty, open, and convex.  Its
linearization is a Banach-space isomorphism
\[
 U:\F(\Omega,o)\longrightarrow\F(\Omega',0),
 \qquad U\delta(x)=\delta(Ax),
\]
with
\[
 \|U\|\leq\|T\|,
 \qquad \|U^{-1}\|\leq\|T^{-1}\|.
\]
Lemma~\ref{lem:euclidean} gives a norm-one map
$Q_2:\F(\Omega',0)^{**}\to\F(\Omega',0)$ that is the identity on the
canonical copy.  Define
\[
 P_\Omega=U^{-1}Q_2U^{**}.
\]
Then
\[
 \|P_\Omega\|
 \leq\|U^{-1}\|\,\|Q_2\|\,\|U^{**}\|
 \leq\|T^{-1}\|\,\|T\|
 =d_{BM}(E,\ell_2^d).
\]
For $u\in\F(\Omega,o)$, naturality gives
\[
 P_\Omega\kappa_{\F(\Omega,o)}u
 =U^{-1}Q_2\kappa_{\F(\Omega',0)}Uu
 =U^{-1}Uu=u.
\]
This is the required projection.
\end{proof}

\begin{remark}[The one-dimensional case]
The source theorem and Question 1.5 are framed around $d\geq2$.  The same
argument also covers $d=1$: the source paper recalls that $\F(\R)$ is
isometric to $L^1(\R)$ and is $1$-complemented in its bidual, while closed
intervals are nonexpansive retracts of $\R$.
\end{remark}

\section{Verification notes and scope}

The proof was audited at the following possible failure points.
\begin{enumerate}
\item \textbf{Open versus closed domain.}  No projection onto the open set is
claimed.  The nearest-point map lands in $\overline\Omega$, and
Lemma~\ref{lem:completion} proves that this has exactly the same free space as
$\Omega$.
\item \textbf{Metric projection.}  Nonexpansiveness is used only after moving
to a Euclidean norm.  It need not hold for the original norm.
\item \textbf{Basepoint.}  The Euclidean retraction fixes every point of the
closed convex set, in particular the chosen basepoint.  The affine map $A$
sends that basepoint to zero before linearization.
\item \textbf{Norm accounting.}  The Euclidean stage has norm one.  The only
loss is conjugation by $U$, bounded by
$\|T\|\|T^{-1}\|=d_{BM}(E,\ell_2^d)$.
\item \textbf{No circularity.}  The only result imported from the source
paper is its already proved whole-space Euclidean case (Corollary 1.3), not
the open-domain assertion under question.
\end{enumerate}

This packet answers Question 1.5 only.  It does not settle the universal
dimension-free constant in Question 1.4 or complementability of
$\F(\ell^2)$ in Question 1.6.

\section{Novelty-search bounds and review recommendation}

The run's cheap indexes were searched for arXiv:1703.08384 and the core terms
``open convex'', ``Lipschitz-free'', ``complemented in its bidual'', and
``Banach--Mazur distance''.  Bounded arXiv/web searches on 9 August 2026 used
the exact Question 1.5 sentence, the exact source title, and close variants of
those terms.  They recovered the source paper and related later work, but no
direct later proof of this statement.  This is evidence only within the
searched corpus, not a claim of exhaustive novelty.

Specialist review should focus on the functorial sandwich
$RQ_HI^{**}$ and on whether a later publication already records this short
retraction argument.  Mathematically, the key observations are standard and
all norm estimates are explicit.

\begin{thebibliography}{9}
\bibitem{CKK}
M.~C\'uth, O.~F.~K. Kalenda, and P.~Kaplick\'y,
\emph{Finitely additive measures and complementability of Lipschitz-free
spaces}, arXiv:1703.08384v2 (2018).

\bibitem{McShane}
E.~J. McShane,
\emph{Extension of range of functions},
Bulletin of the American Mathematical Society \textbf{40} (1934), 837--842.
\end{thebibliography}

\end{document}
